\chapter{Introducción}
\label{cap:intro}

% GUION. Fuentes: doc/memoria/estructura.md, doc/00_Justificaciones.md, doc/03_Ruido_Cuantico.md
% EXTENSION: 3-5 paginas (Instrucciones seccion 2.3). No mas.
%
% ⚠️ ESTRUCTURA FIJADA POR LA NORMA (seccion 2.3): Motivacion · Planteamiento del
%    problema · Estructura del trabajo. Aqui los objetivos van «a grandes rasgos»; los
%    objetivos formales, SMART, viven en el capitulo 5.
%
% ⚠️ DECISION DE REDACCION (sep-2026): la introduccion ANTICIPA el hallazgo negativo.
%    Motivo: §1.3 ya lo enuncia, y plantear §1.2.2 como pregunta abierta dejaba el
%    capitulo contradiciendose consigo mismo veinte lineas despues. Ademas, en un
%    trabajo de tipo 3 el tribunal espera saber desde el principio que se compara y
%    que salio.
%
% ⚠️ REVISAR AL CERRAR RESULTADOS: las contribuciones solo se enuncian del todo bien
%    cuando se sabe que se ha demostrado.

% Parrafo de entrada del capitulo. Solo dice QUE hay en cada apartado; el contenido no se
% adelanta aqui, que para eso estan los apartados.
En este capítulo se introduce el próposito de esta memoria. El apartado~\ref{sec:motivacion} explica la utilidad de un ordenador cuántico y cómo el ruido les perjudica. El
apartado~\ref{sec:planteamiento} plantea el problema concreto que este trabajo aborda: corregir el ruido
exige ejecutar el circuito, y ejecutarlo es caro; y enuncia la propuesta. Finalmente, el
apartado~\ref{sec:estructura} describe la organización de la memoria.

\section{Motivación}
\label{sec:motivacion}
% ⚠️ FUENTES RESTRINGIDAS POR DECISION DEL AUTOR (sep-2026): esta seccion 1.1 se apoya
%    UNICAMENTE en tres fuentes divulgativas:
%      1.1.1 -> ibm2026computacion   (IBM Think, «¿Que es la computacion cuantica?»)
%      1.1.2 -> tqi2026nisq          (The Quantum Insider, guia NISQ)
%      1.1.3 -> duque2025ruido       (Manu Duque, «Ruido cuantico»)
%    NO añadir aqui Shor, Grover, Preskill, Kim, Cai ni Nielsen: se citan mas adelante.


La computación cuántica es un campo de la informática y la ingeniería que aprovecha
las cualidades propias de la mecánica cuántica para resolver problemas que quedan fuera del
alcance incluso de los ordenadores clásicos más potentes \cite{ibm2026computacion}. Y esta definición
hay que leerla con cuidado porque para la mayoría de las tareas se espera que los ordenadores clásicos sigan siendo la mejor
solución. La ventaja cuántica no es general, sino que está
concentrada en un tipo muy concreto de problema.

Ese tipo de problema es aquel en el que muchas variables interactúan de forma complicada. 
Un ejemplo es modelar el
comportamiento de los átomos de una molécula, complejo por las múltiples interacciones entre
sus electrones \cite{ibm2026computacion}. Un ordenador clásico necesita
\emph{calcular} el comportamiento de la molécula, mientras que un ordenador cuántico no
necesita hacer los cálculos, porque puede \textbf{imitar} el sistema molecular.

% ⚠️ FORMATO: titulo arriba y fuente abajo (Instrucciones §1.3, p. 6).
% FOTOGRAFIA DE TERCEROS. Origen: la URL que aporto el autor,
%   https://i0.wp.com/clipset.com/wp-content/uploads/2018/01/ibm-cuantico4.jpg
% ⚠️ SIN LICENCIA DECLARADA en la fuente. Casi con seguridad es material grafico de IBM
%    republicado por el medio, o una foto tomada en la feria. Antes de la entrega hay que
%    localizar el origen primario y su licencia, o sustituirla por una del banco de
%    imagenes de IBM Research, que si publica condiciones de uso.
% ⚠️ Es una JPG de 960x1280 px. A 0.35\textwidth (5,3 cm) da ~460 dpi, de sobra para los
%    300 dpi que exige el README de figuras. NO ampliarla mas de 0.6\textwidth.
\begin{figure}[htbp]
    \centering
    \includegraphics[width=0.35\textwidth]{cap01/ibm_cuantico}

    \caption{Un procesador cuántico de IBM instalado en su criostato.}
    \label{fig:ibm_cuantico}

    {\footnotesize Fuente: Clipset (2018).}
\end{figure}

Lo que se ve en la Figura~\ref{fig:ibm_cuantico} no es el ordenador cuántico, sino el
frigorífico que lo contiene. El procesador cuántico es una
pieza del tamaño de una moneda que va sujeta abajo del todo, a una temperatura cercana al
cero absoluto. Toda esa maquinaria existe para un solo motivo: aislar los qubits del entorno,
porque cualquier vibración o cualquier resto de calor los perturba. Esa fragilidad es
justamente el problema del que trata esta memoria.

En general para hacerse una idea, los ordenadores cuánticos sirven, o se prevé que servirán, para solventar problemas
complejos en campos como: industria farmacéutica, química y ciencia de materiales, energía, finanzas, biología computacional, aprendizaje automático, ciberseguridad y semiconductores \cite{ibm2026computacion}.  Históricamente, la primera aplicación
práctica que se propuso fue el algoritmo de factorización de enteros
publicado por Peter Shor en 1994 \cite{shor1994algorithms}, que mostró cómo una máquina cuántica podría llegar a
romper sistemas criptográficos.

Sin embargo, aunque la computación cuántica sea una tecnología tan prometedora, todavía presenta
ciertos desafíos y limitaciones importantes. 

Actualmente, todos los ordenadores cuánticos disponibles comercialmente (los de IBM, Google, IonQ\dots) pertenecen a una misma categoría, la llamada \textbf{era NISQ}, \textit{Noisy Intermediate-Scale Quantum} \cite{tqi2026nisq}. El físico John Preskill acuñó este término en 2018 para describir una generación de ordenadores cuánticos que ya son difíciles de simular con ordenadores clásicos, pero que siguen teniendo demasiados errores como para aprovechar todo el potencial en aplicaciones reales.
\cite{tqi2026nisq}.

Los qubits, la unidad básica de información en la computación cuántica, son extremadamente sensibles a cualquier perturbación de su entorno, y cada operación introduce una pequeña posibilidad de error. Estos errores se acumulan con el tiempo y limita la fiabilidad de los resultados obtenidos. Lo actuales sistemas disponen de entre decenas y unos pocos miles de qubits que no son suficientes para que lo ordenadores cuánticos sean totalmente fiables. \cite{tqi2026nisq}.

Aunque se conoce una posible solución, llamada Correción de Errores Cuánticos (QEC, por sus siglas en inglés \textit{Quantum Error Correction}), llevarla a la práctica sigue siendo uno de los mayores desafíos de la computación cuántica. Para hacerse una idea: para conseguir una buena fidelidad de un solo qubit es necesario ejecutar cientos o miles de ellos \cite{tqi2026nisq}. Esto hace que los ordenadores cuánticos actuales estén aún lejos de poder aplicarse de manera eficaz.

Mientras tanto, existe una alternativa práctica en la que se centra nuestro trabajo: la Mitigación de Errores Cuánticos (QEM, por sus siglas en inglés \textit{Quantum Error Correction}). En lugar de intentar eliminar el ruido, su objetivo es estimar cuál habría sido el resultado obtenido en un escenario ideal sin errores. Algunas de las técnicas que se incluyen en este área son: \textit{Zero Noise Extrapolation} (ZNE), que estima el resultado ideal extrapolando las mediciones obtenidas de ruido; \textit{Probabilistic Error Cancellation} (PEC), que compensa los efectos del ruido mediante una combinación probabilística; y \textit{Readout Error Mitigation} (REM), que corrige los errores producidos durante la etapa de medición de los qubits..


\section{Planteamiento del problema}
\label{sec:planteamiento}

\subsection{Limitaciones de los enfoques actuales QEM}
\label{sec:limitaciones}

El problema de estos métodos no es la precisión, sino el \textbf{coste}. Tanto ZNE como PEC
necesitan ejecutar el circuito \textbf{varias veces más} para poder estimar cuánto lo ha desviado el ruido. Como consecuencia, 
la corrección de errores necesita una 
cantidad adicional de tiempo de ejecución en el procesador cuántico, un recurso escaso y caro en la actualidad.

De ahí el interés, en este ámbito, del \textbf{aprendizaje automático} (ML-QEM, \textit{Machine Learning - Quantum Error Mitigation}): si un modelo aprende la relación entre el ruido y su efecto, la corrección se calcula sin gastar ejecuciones adicionales. El procedimiento general consta de tres pasos \cite{liao2024machine}:

\begin{enumerate}
      \item \textbf{Generación.} Se produce una gran cantidad de circuitos de entrenamiento
          que imitan la estructura del circuito de interés, de forma que el modelo aprenda
          sobre ejemplos parecidos a aquel que después tendrá que corregir.
      \item \textbf{Codificación.} Cada circuito se pasa por un codificador que extrae sus
          características (el número de compuertas, su tipo o su profundidad, entre otras),
          convirtiendo el circuito en una representación que el modelo pueda consumir.
      \item \textbf{Etiquetado.} Los circuitos se ejecutan en la QPU (Quantum Processing Unit) para obtener sus valores esperados con ruido. La comparación con el valor ideal conocido permite calcular la desviación introducida por el ruido, que se utiliza como variable objetivo del modelo.
\end{enumerate}

Con ese conjunto de datos el modelo aprende a estimar la corrección, y una vez entrenado la
aplica sin necesidad de las ejecuciones extra que hacen falta en ZNE o PEC. Ahora bien, el esquema
tiene una limitación: 

\begin{itemize}
      \item el modelo toma como entrada el \textbf{valor ruidoso ya medido}, y por tanto \textbf{es necesario haber ejecutado el circuito} para poder corregirlo. La corrección sigue llegando después del gasto,
      \item y el \textbf{valor ideal} debe conocerse de antemano para entrenar el modelo.
\end{itemize}

% ⚠️ FORMATO: titulo arriba y fuente abajo (Instrucciones §1.3, p. 6).
% Figura propia. Fuente editable en figuras/cap01/mlqem_flujo.svg; el PDF se regenera con
%   rsvg-convert -f pdf -o mlqem_flujo.pdf mlqem_flujo.svg
% Representa la enumeracion de esta misma seccion Y la limitacion que la sigue: el trazo
% rojo discontinuo es el valor ruidoso medido entrando como entrada del modelo, que es
% exactamente lo que obliga a ejecutar el circuito antes de poder corregirlo. Sin ese
% trazo la figura seria una ilustracion redundante del texto de al lado.
\begin{figure}[htbp]
    \centering
    \includegraphics[width=0.92\textwidth]{cap01/mlqem_flujo}

    \caption{El esquema general de la mitigación por aprendizaje automático.}
    \label{fig:mlqem}

    {\footnotesize Fuente: Claude Code \cite{anthropic_claude}.}
\end{figure}

Para conseguir el valor ideal existen varios métodos pero se ejecutan en ordenadores clásicos. Como mostraremos más adelante, los ordenadores clásicos son incapaces de reproducir los procesos de los ordenadores cuánticos a partir de un número de qubits.

\subsection{La pregunta de este trabajo}
\label{sec:pregunta}

De ahí surge la pregunta que se trata en la memoria:

\begin{center} 
      \doublebox{
      \begin{minipage}{0.85\linewidth} \vspace{0.5em}
             \centering 
             \textit{¿Se puede predecir la degradación que sufrirá un circuito \textbf{antes de ejecutarlo}, mirando únicamente su estructura y la calibración del procesador?} \vspace{0.5em} \end{minipage}} 
\end{center}

Lo que se propone no es cambiar el método, sino el momento en que se aplica la corrección (véase el Capítulo~\ref{sec:reformulacion}): en lugar de corregir el resultado una vez obtenido, intentar \textbf{anticiparlo}. Si esto fuera posible, su utilidad sería inmediata y diferente a la de los métodos anteriores. Permitirá estimar de antemano qué versión de un circuito merece la pena ejecutar, o incluso si merece la pena ejecutarlo, sin consumir tiempo de procesador cuántico. Aunque la literatura ha explorado ideas similares, como clasificar circuitos según la calidad esperada de sus resultados antes de ejecutarlos \cite{hartnett2024learning}, todavía el enfoque `evitando recursos cuánticos' sigue contando con escasas contribuciones. 

\subsection{Propuesta}
\label{sec:propuesta}

Responder a esa pregunta exige, antes que un modelo, un conjunto de datos que no existe.
De modo que lo primero que se propone es construirlo: \textbf{veinte mil circuitos} de entre
cinco y quince qubits, simulados sobre el gemelo digital de un procesador de IBM y
etiquetados con la degradación \textbf{exacta} que sufre cada uno. El ruido no se inventa, se
reconstruye día a día a partir de la \textbf{calibración real} del procesador durante cuarenta y
dos días. Y la partición se diseña para que el conjunto de prueba obligue a
generalizar en \textbf{tres ejes independientes} (el tipo de circuito, su tamaño y el día de
calibración).

Sobre esos datos se comparan \textbf{tres soluciones} bajo un mismo protocolo: una regresión
lineal regularizada (\textbf{Ridge}), un \textbf{Random Forest} (el mejor modelo tabular de referencia
según la ablación de \citeA{liao2024machine}) y una \textbf{red neuronal de grafos con paso de
mensajes} (MPNN, \textit{message passing neural network}). La elección no es arbitraria,
porque un circuito no es una fila: es una matriz de altura variable, una fila por puerta.
Los dos primeros modelos exigen filas de longitud fija, así que hay que \textbf{resumir} el
circuito en columnas, y ese resumen tira el orden y la estructura (dos circuitos
con las mismas puertas estructuradas de formas distintas generan la misma fila). Sin embargo, la red de grafos consume el circuito tal cual. La
diferencia entre unos y otra es, por tanto, \textbf{la medida de cuánto vale la estructura que el
resumen descarta}. Este planteamiento está respaldado por trabajos previos. Por ejemplo, \citeA{tudisco2025graph} demuestra, antes de introducir su red basada en grafos, que características globales como el número de qubits, la profundidad del circuito o el recuento de puertas no son suficientes para distinguir adecuadamente las clases consideradas. Además, la comparación entre redes de grafos y Random Forest aplicada a circuitos cuánticos ya cuenta con un antecedente en \citeA{martyniuk2025benchmarking}.

La última decisión de diseño que distingue esta propuesta de los enfoques habituales en la literatura es que los modelos no intentan predecir un único valor, sino cinco observables al mismo tiempo. La decisión de escoger varias en lugar de una se debe a que existen tipos de circuitos donde varios observables valen cero y no hay nada que predecir. Ampliando el número de observables a cinco, conseguimos que todas las familias de circutios aporten señales para el entrenamiento del modelo.


% =====================================================================================
% ⚠️ SECCION 1.3 «CONTRIBUCIONES» RETIRADA (sep-2026) POR DECISION DEL AUTOR.
%
% Se conserva comentada porque el material es reutilizable: las contribuciones 1 y 3
% estan enunciadas en §1.2.3, y la 2 y la 4 son los capitulos 8 y 7. Si se decide
% reponer la seccion, basta con descomentar el bloque de abajo y REVISARLO contra los
% resultados, que era su aviso original.
% =====================================================================================
% \section{Contribuciones}
% \label{sec:contribuciones}
% % Fuente: doc/memoria/materiales.md y doc/00_Justificaciones.md §4.
% % ⚠️ REVISAR AL CERRAR RESULTADOS.
% % ⚠️ Deliberadamente SIN CIFRAS de R²: las de materiales.md §... son del dataset v1 y la
% %    regla del proyecto prohibe citarlas. Reponer las del v2 cuando esten medidas.
%
% \begin{enumerate}
%     \item \textbf{Un conjunto de datos que no existía.} Veinte mil circuitos etiquetados con
%           su degradación exacta, generados sobre la calibración real de un procesador de IBM
%           a lo largo de $42$ días, y organizados en tres ejes de generalización separables
%           (tipo de circuito, tamaño y tiempo) que permiten medir por separado tres formas
%           distintas de fallar fuera de distribución.
%
%     \item \textbf{La demostración medida de que el objetivo natural del problema no es
%           predecible antes de ejecutar}, junto con su explicación: la degradación es el
%           producto de un factor que sí depende del circuito y del hardware por otro que
%           depende del algoritmo y que es, precisamente, la incógnita para la que haría falta
%           el ordenador cuántico. De ahí se sigue una reformulación del objetivo que sí es
%           aprendible, y que es la que sostiene el resto del trabajo.
%
%     \item \textbf{Una comparativa controlada de tres modelos que mide si la representación
%           en grafo aporta información real.} Ridge, Random Forest y una red neuronal de
%           grafos resuelven el mismo problema con el mismo protocolo y la misma partición; la
%           diferencia entre los dos primeros y el tercero es exactamente el valor de la
%           estructura que la agregación tabular descarta.
%
%     \item \textbf{La detección, medición y corrección de un error en las propias
%           etiquetas}, junto con la razón metodológica por la que la comprobación original no
%           lo detectó y el procedimiento con el que se validó la regeneración.
% \end{enumerate}

\section{Estructura de la memoria}
\label{sec:estructura}
% ⚠️ REESCRITA (sep-2026). La version anterior remitia a \ref{cap:fundamentos} y
%    \ref{cap:sota} como si fueran CAPITULOS, y desde la reestructuracion son las
%    SECCIONES 2.1 y 2.2: el texto imprimia «los capitulos 2.1 y 2.2». Corregido.
%
% 🔴 SI SE MUEVE, FUNDE O PARTE UN CAPITULO, HAY QUE REHACER ESTA LISTA. Es el unico
%    sitio de la memoria donde se enumeran los siete de corrido, asi que un desajuste
%    aqui se ve a la primera.
%
% Guion acordado: una lista, un parrafo corto por capitulo, y que el numero de capitulos
% se vea de un vistazo. Nada de prosa corrida.

La memoria se organiza en \textbf{siete capítulos} y dos anexos:

\begin{enumerate}
    \item \textbf{Introducción.} Presenta para qué sirve un ordenador cuántico y por qué el
          ruido limita hoy lo que puede resolver, delimita el problema concreto que se aborda y
          enuncia la propuesta.

    \item \textbf{Contexto y estado del arte.} Reúne lo imprescindible para leer el resto —del
          qubit a las puertas y al circuito como grafo acíclico dirigido— y revisa lo publicado,
          ordenado por cuándo actúa cada familia de métodos respecto a la ejecución. Cierra
          marcando el hueco que este trabajo ocupa.

    \item \textbf{Objetivos y metodología de trabajo.} Fija el objetivo general y los cinco
          específicos, la métrica que decide entre modelos y el proceso con el que se llega a
          ellos.

    \item \textbf{Planteamiento de la comparativa.} Es el capítulo de diseño. Establece qué se
          predice (los observables, el factor de supervivencia y el umbral de señal), de dónde
          salen los datos y cómo se parten en tres ejes de generalización, y con qué métricas se
          juzga.

    \item \textbf{Desarrollo de la comparativa.} Ataca los tres modelos uno a uno: qué datos ve
          cada uno, qué preprocesado es exclusivamente suyo y con qué arquitectura e
          hiperparámetros se quedó.

    \item \textbf{Discusión y análisis de resultados.} Presenta lo que salió, separando dos
          preguntas distintas: cuánto acierta cada modelo al predecir el factor, y cuánto mejora
          el observable una vez aplicado.

    \item \textbf{Conclusiones y trabajo futuro.} Enlaza cada objetivo con su conclusión, declara
          las limitaciones del trabajo y propone las líneas que las cierran.
\end{enumerate}

Los anexos recogen las herramientas empleadas y la declaración de uso de herramientas de
inteligencia artificial.
